% ======================================================================
% Chapter 2 -- Part C: criteria, trade-off table, failure walk-through,
% decision record, review checklist, end matter
% ======================================================================

\section{Eight Evaluation Criteria}
\label{sec:criteria}

To compare designs we need fixed yardsticks. Every case in the book uses the same eight criteria, so that designs can be compared across cases and so that no one quietly drops the criterion on which their favorite design looks worst. Table~\ref{tab:criteria} defines each criterion and anchors its scale: what a score of 1 and a score of 5 look like, and how the score is eventually measured rather than estimated.

\begin{table}[htbp]
\centering
\small
\caption{The eight evaluation criteria, with scale anchors and measurement.}
\label{tab:criteria}
\begin{tabularx}{\textwidth}{@{}>{\raggedright\arraybackslash}p{2.5cm}>{\raggedright\arraybackslash}X>{\raggedright\arraybackslash}X>{\raggedright\arraybackslash}p{3.1cm}@{}}
\toprule
Criterion & 1 (poor) & 5 (excellent) & Measured by \\
\midrule
Quality & Frequent wrong outcomes, some unnoticed & Outcomes at or above the best current human level & Task accuracy on an evaluation set; KPI error rates \\
Cost & Run cost per transaction erodes the value case & Run cost negligible next to value per transaction & Cost per transaction (model, infrastructure, tool calls) \\
Speed & No faster than today & Meets the speed KPI with margin & End-to-end latency or cycle time, median and 95th percentile \\
Reliability & Outages or erratic behavior under load or bad input & Degrades gracefully; behavior consistent run to run & Success rate, variance across repeated runs, fallback rate \\
Security and privacy & Broad permissions; untrusted input reaches agents with write tools & Least privilege; trust boundaries guarded & Permission review; red-team results \\
Control and auditability & Cannot explain why an outcome happened & Every outcome traceable to inputs, rules, sources and approvers & Fraction of outcomes with a complete trace \\
Maintainability & Behavior changes in surprising places when anything changes & Rules, prompts, tools and models change independently & Effort to make a standard change; regression rate \\
Human workload & Adds work or moves it to a scarcer role & Removes routine work; people handle only real exceptions & Minutes per transaction by role; queue sizes \\
\bottomrule
\end{tabularx}
\end{table}

Three remarks on using the criteria.

First, \emph{at design time most scores are estimates}. You score from experience, from the pattern cards, and from quick prototypes. The ``measured by'' column says how each estimate will be replaced by a number in the build and evaluation sections of each case. A trade-off table whose scores were never checked against measurements is an opinion with a grid around it.

Second, \emph{the criteria are not independent}. Adding a checkpoint raises quality and control and lowers speed and, often, human workload. That coupling is the reason trade-offs exist at all.

Third, \emph{some requirements are not trade-offs}. A legal or policy requirement---finance approves credits above \$2{,}000; customer data stays in the region---is a constraint. A design that violates a constraint is out, however well it scores. Check constraints first, then score what survives.

% ----------------------------------------------------------------------
\section{The Trade-off Table}
\label{sec:tradeoff}

\begin{definition}
A \keyterm{trade-off table} scores two or more alternative designs on the eight criteria, combines the scores with weights that reflect the business context, and records how sensitive the result is to those weights.
\end{definition}

\subsection{Three designs for returns}

Figure~\ref{fig:returns-designs} shows three designs for Harbor's returns process. All three meet the two policy constraints (inspector judges misuse; finance approves large credits). They differ in how much the system decides on its own.

\begin{description}[leftmargin=1em,itemsep=2pt]
  \item[D1, single returns agent.] One Level~4 agent with tools on all three servers reads the request, applies the warranty rules described in its instructions, and creates the RMA and label itself. This is the do-everything agent from Section~\ref{sec:patterns}. Most teams build it first, and a sponsor will always ask about it: ``why not just one agent?''
  \item[D2, rep-reviewed workflow.] A Level~3 router classifies the request; code checks warranty eligibility against the table; a Level~1 agent drafts the RMA and a customer reply; a representative approves every RMA before it is created.
  \item[D3, straight-through with exceptions.] As D2, but the representative's approval is required only when the case falls outside a safe zone: code issues the RMA automatically when the item is in warranty, the value is under \$500, the router's confidence is high, and the customer is not flagged by a repeat-return rule. Everything else goes to a representative.
\end{description}

\begin{figure}[htbp]
\centering
\begin{tikzpicture}[x=1cm, y=1cm, sg base/.append style={text width=1.7cm, minimum height=7.5mm, font=\scriptsize, inner sep=2pt}]
% ---- D1
\node[font=\small\bfseries, anchor=west] at (-1,0.9) {D1};
\node[sg trigger] (c1) at (0,0) {Customer};
\node[sg agent] (a1) at (2.6,0) {Returns agent\\{\tiny L4; warranty in prompt}};
\node[sg mcp, text width=2.4cm] (m1) at (5.9,0) {ERP, WMS, Ship\nodepart{two}{\tiny\ttfamily 9 tools incl.\ create\_rma}};
\node[sg human] (f1) at (9.1,0) {Finance\\{\tiny $>$\$2k credits}};
\draw[sg flow] (c1) -- (a1); \draw[sg call] (a1) -- (m1); \draw[sg flow] (m1) -- (f1);
% ---- D2
\begin{scope}[yshift=-2.4cm]
\node[font=\small\bfseries, anchor=west] at (-1,0.9) {D2};
\node[sg trigger] (c2) at (0,0) {Customer};
\node[sg agent] (r2) at (2.3,0) {Router\\{\tiny L3}};
\node[sg code] (w2) at (4.6,0) {Warranty\\check};
\node[sg agent] (d2) at (6.9,0) {Draft RMA\\{\tiny L1}};
\node[sg human] (h2) at (9.2,0) {Rep approves\\every RMA};
\node[sg mcp, text width=1.9cm] (m2) at (11.6,0) {ERP, Ship\nodepart{two}{\tiny\ttfamily create\_rma, label}};
\draw[sg flow] (c2) -- (r2); \draw[sg flow] (r2) -- (w2); \draw[sg flow] (w2) -- (d2);
\draw[sg flow] (d2) -- (h2); \draw[sg call] (h2) -- (m2);
\end{scope}
% ---- D3
\begin{scope}[yshift=-5.0cm]
\node[font=\small\bfseries, anchor=west] at (-1,0.9) {D3};
\node[sg trigger] (c3) at (0,0) {Customer};
\node[sg agent] (r3) at (2.3,0) {Router\\{\tiny L3}};
\node[sg code] (w3) at (4.6,0) {Warranty \&\\safe-zone rules};
\node[sg agent] (d3) at (6.9,0) {Draft RMA\\{\tiny L1}};
\node[sg human] (h3) at (9.2,-1.1) {Rep reviews\\exceptions};
\node[sg mcp, text width=1.9cm] (m3) at (11.6,0) {ERP, Ship\nodepart{two}{\tiny\ttfamily create\_rma, label}};
\draw[sg flow] (c3) -- (r3); \draw[sg flow] (r3) -- (w3); \draw[sg flow] (w3) -- (d3);
\draw[sg call] (d3) -- node[sg lab, above] {[safe zone]} (m3);
\draw[sg flow] (w3.south) |- node[sg lab, pos=0.75, below] {[else]} (h3.west);
\draw[sg call] (h3.east) -| node[sg lab, pos=0.25, below] {[approved]} (m3.south);
\end{scope}
\end{tikzpicture}
\caption{Three designs for returns, simplified to the intake and RMA stage. The inspection, event trigger and credit stages, which are the same in D2 and D3, are omitted.}
\label{fig:returns-designs}
\end{figure}

\subsection{Scoring and weighting}

Table~\ref{tab:tradeoff} scores the three designs. Each score has a reason, and the reasons matter more than the numbers. D1 scores 1 on control because no one can show afterwards which warranty rule it applied; it scores 5 on speed and workload because it never waits for anyone. D2 is the mirror image. D3 gives up a point of quality and a point of control relative to D2 to gain two points on speed and workload.

The weights come from the business context, not from the designer's taste. Harbor's sponsor cares most about speed and correct credits; the controller adds weight to auditability; maintainability matters least because the process changes rarely. Writing the weights down before scoring keeps them honest, and showing them to the sponsor turns a technical argument into a business conversation.

The weighted score of design $d$ is
\begin{equation}
  S_d = \sum_{c=1}^{8} w_c \, s_{d,c}, \qquad \sum_{c=1}^{8} w_c = 1,
  \label{eq:weighted}
\end{equation}
where $s_{d,c}$ is design $d$'s score on criterion $c$ and $w_c$ is the criterion's weight.

\begin{table}[htbp]
\centering
\small
\caption{Trade-off table for the returns designs. Scores 1--5; weights from the business context.}
\label{tab:tradeoff}
\begin{tabular}{@{}lcccc@{}}
\toprule
Criterion & Weight & D1 single agent & D2 rep-reviewed & D3 straight-through \\
\midrule
Quality & 0.20 & 2 & 5 & 4 \\
Cost & 0.10 & 3 & 4 & 4 \\
Speed & 0.20 & 5 & 2 & 4 \\
Reliability & 0.10 & 2 & 4 & 4 \\
Security and privacy & 0.10 & 2 & 4 & 4 \\
Control and auditability & 0.15 & 1 & 5 & 4 \\
Maintainability & 0.05 & 3 & 4 & 3 \\
Human workload & 0.10 & 5 & 2 & 4 \\
\midrule
Weighted score $S_d$ & 1.00 & 2.90 & 3.75 & \textbf{3.95} \\
\bottomrule
\end{tabular}
\end{table}

\subsection{Is the decision robust?}

D3 leads, but by only 0.20 points on a five-point scale, and every score in the table is an estimate. Before recording a decision, ask how wrong the weights would have to be for the answer to change. Listing~\ref{lst:tradeoff} does this in two ways. It draws ten thousand weight vectors scattered around the chosen weights and counts how often each design wins; and for each criterion, it finds how far that one weight would have to rise (with the others scaled down in proportion) before D2 overtakes D3.

\lstinputlisting[language=Python, caption={Weighted scores and weight sensitivity for Table~\ref{tab:tradeoff} (\texttt{ch02/tradeoff.py}).}, label={lst:tradeoff}, firstline=2]{\chtwocode tradeoff.py}

Running it prints:
\begin{lstlisting}[language={}, numbers=none, frame=none, xleftmargin=1.5em]
D1 single agent        2.90
D2 rep-reviewed        3.75
D3 straight-through    3.95
D1 single agent        wins 0.0% of perturbed weightings
D2 rep-reviewed        wins 15.1% of perturbed weightings
D3 straight-through    wins 84.9% of perturbed weightings
D2 overtakes D3 if 'quality' weight rises from 0.20 to 0.34
D2 overtakes D3 if 'control' weight rises from 0.15 to 0.30
D2 overtakes D3 if 'maintainability' weight rises from 0.05 to 0.21
\end{lstlisting}

Three conclusions follow. D1 never wins under any plausible weighting, so the answer to ``why not one agent?'' is not a matter of taste. D3 wins in about 85\% of the perturbed weightings, so the choice is reasonably robust. And the result names exactly what would reverse it: if Harbor's controller valued auditability twice as much as assumed, or if quality mattered substantially more than speed, D2 would be the better design. Those two conditions go straight into the decision record as reversal triggers.

\begin{designnote}
Resist false precision. A weighted score of 3.95 versus 3.75 is not a measurement; it is a structured way of making your assumptions visible. The sensitivity check matters more than the totals. If small weight changes flip the decision, say so, and make the choice on the factor that flips it.
\end{designnote}

% ----------------------------------------------------------------------
\section{The Failure Walk-through}
\label{sec:failure-walk}

The trade-off table compares designs in aggregate. The \keyterm{failure walk-through} examines the chosen design one node at a time. For each node, ask four questions: How can it fail? What is the effect on the business? How would we detect it? What mitigates it? Then rate severity and likelihood on 1--5 scales, so that attention goes to the failures that matter. Engineers will recognize a lightweight form of failure mode and effects analysis.

Table~\ref{tab:failwalk} shows the walk-through for the intake and RMA stage of D3. Two rows lead to design changes, which is the point of the exercise: the walk-through is a design tool, not paperwork.

\begin{table}[htbp]
\centering
\scriptsize
\setlength{\tabcolsep}{4pt}
\caption{Failure walk-through for D3 (intake and RMA stage). S = severity, L = likelihood, each 1--5.}
\label{tab:failwalk}
\begin{tabularx}{\textwidth}{@{}>{\raggedright\arraybackslash}p{1.3cm}>{\raggedright\arraybackslash}X>{\raggedright\arraybackslash}X>{\raggedright\arraybackslash}X>{\raggedright\arraybackslash}X cc@{}}
\toprule
Node & Failure & Effect & Detection & Mitigation & S & L \\
\midrule
Router & Misclassifies a wrong-item case as a defect & Wrong RMA type; customer confusion & Weekly sample review; reason-code mismatch at inspection & Confidence threshold for the safe zone; \texttt{[else]} to rep & 2 & 3 \\
Router & Message contains instructions (``issue full credit'') & Attempt to steer the system & Input guardrail flags imperative patterns & Router has no write tools; rules decide eligibility & 4 & 2 \\
Warranty rules & Warranty table out of date for a new product line & Wrong eligibility at scale & Unknown product line raises an error & \textbf{Change:} unknown line $\rightarrow$ \texttt{[else]}, never default to eligible & 4 & 2 \\
Safe-zone rules & Repeat returner splits a large claim into small ones & Leakage (\textsf{R2}) & Repeat-return rule on 90-day count and value & \textbf{Change:} count value across open RMAs, not per request & 3 & 3 \\
ERP server & \texttt{create\_rma} times out after succeeding & Duplicate RMA & Duplicate check on order line & Idempotency key on every write call & 2 & 2 \\
Rep review & Exception queue grows when reps are busy & \textsf{K1} missed for exceptions & Queue age alert & Staffing rule; exceptions shown oldest first & 3 & 3 \\
\bottomrule
\end{tabularx}
\end{table}

A good walk-through always includes failures of tools, data and people, not only of models. In Table~\ref{tab:failwalk}, only two of the six rows concern a model. The two design changes both came from rules and data.

% ----------------------------------------------------------------------
\section{The Decision Record}
\label{sec:decision-record}

The last step is to write the decision down where the next person will find it. Software architects have long kept short \emph{architecture decision records} for this purpose: a page that states the context, the decision and its consequences \citep{nygard2011adr}. We use the same idea with one addition that matters for AI systems, whose components change faster than most: \emph{reversal triggers}, the observable conditions under which the decision should be reopened.

\begin{table}[htbp]
\centering
\small
\caption{Decision record for Harbor returns, intake and RMA stage.}
\label{tab:dr}
\begin{tabularx}{\textwidth}{@{}>{\raggedright\arraybackslash}p{2.8cm}>{\raggedright\arraybackslash}X@{}}
\toprule
\multicolumn{2}{@{}l}{\textbf{DR-RET-001: Issue low-risk RMAs automatically; route exceptions to a representative}} \\
\midrule
Status and date & Accepted, pilot phase. Owner: returns manager. \\
Context & 1{,}100 returns a month; median 3 days to RMA; 6\% of credits out of warranty. KPIs \textsf{K1}--\textsf{K3}; risks \textsf{R1}, \textsf{R2}. Constraints: inspector judges misuse; finance approves credits over \$2{,}000. \\
Options considered & D1 single returns agent; D2 rep-reviewed workflow; D3 straight-through with exceptions (Figure~\ref{fig:returns-designs}). \\
Decision & D3, with the two changes from the failure walk-through. Safe zone: in warranty, under \$500, router confidence at or above the evaluated threshold, customer not flagged. \\
Reasons & Highest weighted score (3.95 vs.\ 3.75 and 2.90); wins in about 85\% of perturbed weightings. D1 rejected: no auditable warranty decision. \\
Consequences & Reps handle only exceptions (expected 30--40\% of volume). A weekly sample of automatic RMAs is reviewed. Warranty table becomes a governed configuration with an owner. \\
Reversal triggers & Reopen if: out-of-warranty RMAs in the safe zone exceed 1\% in any month; the controller's weight on auditability doubles (for example, after an audit finding); the exception share exceeds 60\%, which would erase most of the gain over D2. \\
\bottomrule
\end{tabularx}
\end{table}

Note how much of the record came straight from earlier steps: the KPIs and constraints from the method, the options from the patterns, the reasons and two of the reversal triggers from the sensitivity check, and the changes from the walk-through. That is by design. If the record is hard to write, an earlier step was skipped.

% ----------------------------------------------------------------------
\section{Putting the Toolkit Together}
\label{sec:toolkit-together}

Every case chapter uses the toolkit in the same order. As a reminder, and as a checklist for design reviews:

\begin{enumerate}[itemsep=1pt]
  \item Run the seven-step method; number the KPIs and risks.
  \item Draw a context view, then at least two alternative solution graphs. In Part~II, the first alternative is the Google sample as it stands.
  \item Check each graph against the well-formedness rules W1--W8.
  \item Check the policy and legal constraints; drop any design that violates one.
  \item Write the weights from the business context, then score the survivors on the eight criteria, with a reason for each score.
  \item Run the sensitivity check; note what would flip the decision.
  \item Walk the chosen design node by node; fold the resulting changes back into the graph.
  \item Write the decision record, including reversal triggers.
  \item Write node cards for every agent, MCP server and guardrail.
\end{enumerate}

Chapter~\ref{ch:business} adds the business layer that this chapter has taken as given: where KPIs come from, how to estimate the value of moving them, and how to put the cost of errors into numbers, so that the quality and workload scores in a trade-off table can be replaced by dollars.

% ----------------------------------------------------------------------
\section*{Chapter Summary}
\addcontentsline{toc}{section}{Chapter Summary}

\begin{itemize}[itemsep=2pt]
  \item A fixed notation makes the kind, owner and justification of every component visible, makes omissions stand out, and lets designs be compared. The notation has nine node types, five edge types, ownership and trust boundaries, and trace tags.
  \item Eight well-formedness rules (W1--W8) catch the commonest structural mistakes before any discussion of merit.
  \item The seven-step method works from the business inward: goal and metric, actors and systems, judgment versus rules, knowledge, actions, human checkpoints, failure points.
  \item Two patterns appear in almost every case: retrieve-then-answer and the approval gate. Other patterns are introduced when a case first needs them. Three anti-patterns name the commonest failures: the do-everything agent, policy in the prompt and the silent write.
  \item Alternatives are scored on eight fixed criteria with weights from the business context. The sensitivity check matters more than the totals.
  \item The failure walk-through examines tools, data and people as well as models, and usually changes the design.
  \item A decision record captures context, options, decision, reasons, consequences and reversal triggers.
\end{itemize}

\section*{Key Terms}
\addcontentsline{toc}{section}{Key Terms}

\begin{multicols}{2}
\raggedright\small
anti-pattern \\
constraint (vs.\ criterion) \\
decision record \\
design pattern \\
evaluation criteria \\
failure walk-through \\
handoff \\
node card \\
ownership boundary \\
reversal trigger \\
sensitivity check \\
solution graph \\
tool call \\
trace tag \\
trade-off table \\
trust boundary \\
well-formedness rules
\end{multicols}

\section*{Review Questions}
\addcontentsline{toc}{section}{Review Questions}

\begin{enumerate}[itemsep=2pt]
  \item Why does the notation draw the system of record separately from the MCP server in front of it? Which well-formedness rule depends on this?
  \item Explain the difference between a tool call and a handoff. Give an example where choosing the wrong one would leave a case with no owner.
  \item Which well-formedness rules does Team A's design from Chapter~\ref{ch:foundations} violate? List them with one sentence each.
  \item In the returns case, why is warranty eligibility treated as a rule even though representatives found it hard?
  \item When is retrieve-then-answer the wrong pattern? Give an example where a database lookup is the better choice.
  \item When is an approval gate the wrong choice? What would you use instead?
  \item What is the difference between a constraint and a criterion? Why are constraints checked first?
  \item The sensitivity check found that D2 overtakes D3 if the weight on control rises from 0.15 to 0.30. What business event might cause such a change?
  \item Why should a failure walk-through include tools, data and people, not only models? Give an example of each from Table~\ref{tab:failwalk}.
  \item What is a reversal trigger, and why are reversal triggers especially important for AI systems?
\end{enumerate}

\section*{Exercises}
\addcontentsline{toc}{section}{Exercises}

\paragraph{Design track}
\begin{enumerate}[itemsep=2pt]
  \item \textbf{Complete the returns graph.} Extend D3 to include the inspection and credit stages: an event when the warehouse records an inspection, the credit calculation, and the finance gate. Check the full graph against W1--W8.
  \item \textbf{Node cards.} Write node cards, in the format of Table~\ref{tab:nodecard}, for the router and the ERP server in D3.
  \item \textbf{Re-weight.} Suppose Harbor is preparing for an acquisition and its auditors require every RMA to have a named human approver. Is this a constraint or a change of weights? Redo the decision and update the decision record.
  \item \textbf{Pattern hunt.} Pick a published agent sample from any vendor. Draw it in the notation, name the patterns it uses, and list any anti-patterns.
  \item \textbf{Your own process.} Apply the seven-step method to the process you decomposed in Chapter~\ref{ch:foundations}'s exercises, draw two alternatives, and complete a trade-off table with a sensitivity check.
\end{enumerate}

\paragraph{Build track}
\begin{enumerate}[resume, itemsep=2pt]
  \item \textbf{Sensitivity, extended.} Modify Listing~\ref{lst:tradeoff} to vary the \emph{scores} as well as the weights (for example, each score $\pm1$ with some probability). Does D3 remain the choice?
  \item \textbf{Diagram as code.} Using the TikZ styles for the notation in the companion repository, or the Mermaid equivalents provided there, draw D2 and D3 so that they can be versioned alongside the code.
\end{enumerate}

% ----------------------------------------------------------------------
\section*{Lab 2: Measuring Two Designs}
\addcontentsline{toc}{section}{Lab 2: Measuring Two Designs}

\paragraph{Goal.} Replace three of the estimated scores in Table~\ref{tab:tradeoff}---quality, cost and human workload---with measurements, by running the intake stage of D2 and D3 on the same sample of returns.

\paragraph{Time.} About 90 minutes.

\paragraph{Steps.}
\begin{enumerate}[itemsep=2pt]
  \item Open \texttt{ch02/} in the companion repository and install the pinned requirements. The folder contains 200 synthetic return requests with labeled correct outcomes, a warranty table, and stub MCP servers for the ERP and shipping systems.
  \item Run the provided router agent on all 200 requests with \texttt{python run\_router.py}. Record routing accuracy and the confidence score for each request.
  \item Implement the safe-zone rule from the decision record as a plain Python function in \texttt{safe\_zone.py}, including the two changes from the failure walk-through. Write unit tests for it; the router's model must not be involved in these tests.
  \item Run \texttt{python simulate.py --design D2} and \texttt{--design D3}. The script assumes a representative takes 4 minutes per review and reports, for each design, the share of correct outcomes, the model cost per request, and representative minutes per 100 requests.
  \item Choose the router-confidence threshold for the safe zone by sweeping it from 0.5 to 0.95 and plotting correct outcomes against the share of requests handled automatically.
  \item Convert your measurements to 1--5 scores using the anchors in Table~\ref{tab:criteria}, update \texttt{tradeoff.py}, and rerun the sensitivity check.
\end{enumerate}

\paragraph{Deliverable.} An updated trade-off table with measured scores marked, the threshold plot, and a revised decision record. State whether the measurements changed the decision, and if so, which reversal trigger fired.

% ----------------------------------------------------------------------
\section*{Further Reading}
\addcontentsline{toc}{section}{Further Reading}

\begin{itemize}[itemsep=2pt]
  \item \citet{anthropic2024effective} describes the workflow and agent patterns that this chapter and later cases build on.
  \item \citet{nygard2011adr} is the short essay that popularized architecture decision records.
  \item \citet{madaan2023selfrefine} and \citet{shinn2023reflexion} study reflection-style loops, useful background for the reflection pattern in Chapter~\ref{ch:case-fraud}.
\end{itemize}
